\documentclass[11pt]{article}

\usepackage[a4paper,margin=0.75in]{geometry}
\usepackage{amsmath,amssymb}
\usepackage{graphicx}
\usepackage{booktabs}
\usepackage{array}
\usepackage{float}
\usepackage{hyperref}
\usepackage{enumitem}

\setlist{nosep}
\renewcommand{\arraystretch}{1.3}

\begin{document}

%%%%%%%%%%%%%%%%%%%%%%%%%%%%%%%%%%%%%%%%%%%%%%%%%%%%%%%%%%%%

\begin{center}

{\Large\textbf{Shiv Nadar University Chennai}}

\vspace{3mm}

{\large\textbf{CS3807 -- Deep Learning Laboratory}}

\vspace{3mm}

{\Large\textbf{Experiment 3}}

\vspace{2mm}

{\Large\textbf{Implementation of Convolutional Neural Networks (CNNs)}}

{\Large\textbf{for Image Classification}}

\end{center}

\vspace{5mm}

\begin{tabular}{|p{4cm}|p{8cm}|p{2cm}|}
\hline
Degree \& Branch &
B.Tech Artificial Intelligence \& Data Science &
Semester V\\
\hline
Subject Code \& Name &
CS3807 -- Deep Learning Laboratory &
AY: 2026--27\\
\hline
Name & R S Kirithic & \\
\hline
Register Number & 24011101082 & \\
\hline
\end{tabular}

%%%%%%%%%%%%%%%%%%%%%%%%%%%%%%%%%%%%%%%%%%%%%%%%%%%%%%%%%%%%

\section*{1. Objective}

To understand the working principle of Convolutional Neural Networks by
implementing convolution, pooling, feature map visualization, and image
classification using TensorFlow/Keras.

%%%%%%%%%%%%%%%%%%%%%%%%%%%%%%%%%%%%%%%%%%%%%%%%%%%%%%%%%%%%

\section*{2. Learning Outcomes}

Students will be able to

\begin{enumerate}
\item Explain the convolution operation.
\item Compute output dimensions using stride and padding.
\item Visualize feature maps.
\item Implement max pooling and average pooling.
\item Calculate CNN parameters.
\item Build a CNN for image classification.
\item Evaluate CNN performance.
\end{enumerate}

%%%%%%%%%%%%%%%%%%%%%%%%%%%%%%%%%%%%%%%%%%%%%%%%%%%%%%%%%%%%

\section*{3. Background Theory}

A Convolutional Neural Network consists of

\begin{itemize}
\item Convolution Layer
\item Activation Layer
\item Pooling Layer
\item Fully Connected Layer
\end{itemize}

The convolution operation is

\[
Y(i,j)=\sum_m\sum_n X(i+m,j+n)K(m,n)
\]

where

\begin{itemize}
\item $X$ : Input Image
\item $K$ : Kernel (Filter)
\item $Y$ : Feature Map
\end{itemize}

The output feature map size is

\[
\frac{N-F+2P}{S}+1
\]

where

\begin{tabular}{ll}
$N$ & Input Size\\
$F$ & Filter Size\\
$P$ & Padding\\
$S$ & Stride
\end{tabular}

\subsection*{3.1 Common Convolution Kernels}

A \textbf{kernel} (or \textbf{filter}) is a small matrix that slides over
an input image to extract important visual features such as edges,
textures, corners, and smooth regions. During convolution, the kernel is
multiplied element-wise with the corresponding image pixels, and the
results are summed to produce a single value in the output feature map.
Different kernels highlight different image characteristics.

\subsubsection*{1. Identity Kernel}

The identity kernel preserves the original image without modifying its
pixel values.

\[
\begin{bmatrix}
0 & 0 & 0\\
0 & 1 & 0\\
0 & 0 & 0
\end{bmatrix}
\]

\textbf{Purpose:}
\begin{itemize}
\item Preserves the original image.
\item Used for understanding the convolution operation.
\end{itemize}

\subsubsection*{2. Sobel-X Kernel (Vertical Edge Detection)}

The Sobel-X kernel detects \textbf{vertical edges} by measuring intensity
changes along the horizontal direction.

\[
\begin{bmatrix}
-1 & 0 & 1\\
-2 & 0 & 2\\
-1 & 0 & 1
\end{bmatrix}
\]

\textbf{Applications:}
\begin{itemize}
\item Object detection
\item Lane detection
\item Face recognition
\end{itemize}

\subsubsection*{3. Sobel-Y Kernel (Horizontal Edge Detection)}

The Sobel-Y kernel detects \textbf{horizontal edges} by measuring
intensity changes along the vertical direction.

\[
\begin{bmatrix}
-1 & -2 & -1\\
0 & 0 & 0\\
1 & 2 & 1
\end{bmatrix}
\]

\textbf{Applications:}
\begin{itemize}
\item Text detection
\item Boundary extraction
\item Medical image analysis
\end{itemize}

\subsubsection*{4. Laplacian Kernel}

The Laplacian kernel detects edges in all directions using the
second-order derivative of image intensity.

\[
\begin{bmatrix}
0 & -1 & 0\\
-1 & 4 & -1\\
0 & -1 & 0
\end{bmatrix}
\]

\textbf{Applications:}
\begin{itemize}
\item Edge enhancement
\item Satellite image analysis
\item Medical imaging
\end{itemize}

\subsubsection*{5. Sharpening Kernel}

This kernel enhances fine details by emphasizing high-frequency
components.

\[
\begin{bmatrix}
0 & -1 & 0\\
-1 & 5 & -1\\
0 & -1 & 0
\end{bmatrix}
\]

\textbf{Applications:}
\begin{itemize}
\item Image enhancement
\item Optical Character Recognition (OCR)
\item Face recognition
\end{itemize}

\subsubsection*{6. Box Blur Kernel}

The Box Blur kernel smooths an image by replacing each pixel with the
average of its neighboring pixels.

\[
\frac{1}{9}
\begin{bmatrix}
1 & 1 & 1\\
1 & 1 & 1\\
1 & 1 & 1
\end{bmatrix}
\]

\textbf{Applications:}
\begin{itemize}
\item Noise removal
\item Image smoothing
\item Image preprocessing
\end{itemize}

\subsubsection*{7. Gaussian Blur Kernel}

Gaussian blur smooths the image while preserving the overall structure
better than a simple average filter.

\[
\frac{1}{16}
\begin{bmatrix}
1 & 2 & 1\\
2 & 4 & 2\\
1 & 2 & 1
\end{bmatrix}
\]

\textbf{Applications:}
\begin{itemize}
\item Noise reduction
\item Computer vision preprocessing
\item Feature extraction
\end{itemize}

\subsubsection*{8. Emboss Kernel}

The Emboss kernel creates a three-dimensional appearance by emphasizing
directional changes in intensity.

\[
\begin{bmatrix}
-2 & -1 & 0\\
-1 & 1 & 1\\
0 & 1 & 2
\end{bmatrix}
\]

\textbf{Applications:}
\begin{itemize}
\item Texture analysis
\item Artistic image effects
\end{itemize}

\subsubsection*{9. Outline Detection Kernel}

This kernel highlights object boundaries by emphasizing regions with
rapid intensity changes.

\[
\begin{bmatrix}
-1 & -1 & -1\\
-1 & 8 & -1\\
-1 & -1 & -1
\end{bmatrix}
\]

\textbf{Applications:}
\begin{itemize}
\item Image segmentation
\item Boundary detection
\item Shape extraction
\end{itemize}

\subsubsection*{10. Motion Blur Kernel}

This kernel simulates motion blur along the diagonal direction.

\[
\frac{1}{5}
\begin{bmatrix}
1 & 0 & 0 & 0 & 0\\
0 & 1 & 0 & 0 & 0\\
0 & 0 & 1 & 0 & 0\\
0 & 0 & 0 & 1 & 0\\
0 & 0 & 0 & 0 & 1
\end{bmatrix}
\]

\textbf{Applications:}
\begin{itemize}
\item Motion simulation
\item Image restoration
\item Video processing
\end{itemize}

\subsubsection*{Summary of Common Kernels}

\begin{center}
\renewcommand{\arraystretch}{1.3}
\begin{tabular}{|p{3cm}|p{2cm}|p{4cm}|p{4cm}|}
\hline
\textbf{Kernel} & \textbf{Size} & \textbf{Purpose} & \textbf{Applications}\\
\hline
Identity & $3\times3$ & Preserve original image & Baseline comparison\\
\hline
Sobel-X & $3\times3$ & Detect vertical edges & Object detection\\
\hline
Sobel-Y & $3\times3$ & Detect horizontal edges & Boundary detection\\
\hline
Laplacian & $3\times3$ & Detect edges in all directions & Medical imaging\\
\hline
Sharpen & $3\times3$ & Enhance details & Image enhancement\\
\hline
Box Blur & $3\times3$ & Smooth image & Noise removal\\
\hline
Gaussian Blur & $3\times3$ & Reduce noise & Computer vision\\
\hline
Emboss & $3\times3$ & 3D effect & Artistic filtering\\
\hline
Outline & $3\times3$ & Boundary extraction & Segmentation\\
\hline
Motion Blur & $5\times5$ & Simulate motion & Video processing\\
\hline
\end{tabular}
\end{center}

%%%%%%%%%%%%%%%%%%%%%%%%%%%%%%%%%%%%%%%%%%%%%%%%%%%%%%%%%%%%

\section*{4. Numerical Example 1: Convolution}

Consider the following input image.

\[
X=
\begin{bmatrix}
1&2&3\\
4&5&6\\
7&8&9
\end{bmatrix}
\]

Kernel

\[
K=
\begin{bmatrix}
1&0\\
0&1
\end{bmatrix}
\]

Output

First position

\[
1(1)+2(0)+4(0)+5(1)=6
\]

Second position

\[
2(1)+3(0)+5(0)+6(1)=8
\]

Third position

\[
4(1)+5(0)+7(0)+8(1)=12
\]

Fourth position

\[
5(1)+6(0)+8(0)+9(1)=14
\]

Feature Map

\[
\begin{bmatrix}
6&8\\
12&14
\end{bmatrix}
\]

%%%%%%%%%%%%%%%%%%%%%%%%%%%%%%%%%%%%%%%%%%%%%%%%%%%%%%%%%%%%

\section*{5. Numerical Example 2: Max Pooling}

Feature map

\[
\begin{bmatrix}
1&5&2&3\\
7&8&1&0\\
4&6&9&5\\
2&3&1&8
\end{bmatrix}
\]

Using a $2\times2$ max pooling filter,

Window 1

\[
\begin{bmatrix}
1&5\\
7&8
\end{bmatrix}
\rightarrow 8
\]

Window 2

\[
\begin{bmatrix}
2&3\\
1&0
\end{bmatrix}
\rightarrow3
\]

Window 3

\[
\begin{bmatrix}
4&6\\
2&3
\end{bmatrix}
\rightarrow6
\]

Window 4

\[
\begin{bmatrix}
9&5\\
1&8
\end{bmatrix}
\rightarrow9
\]

Output

\[
\begin{bmatrix}
8&3\\
6&9
\end{bmatrix}
\]

%%%%%%%%%%%%%%%%%%%%%%%%%%%%%%%%%%%%%%%%%%%%%%%%%%%%%%%%%%%%

\section*{6. Numerical Example 3: Parameter Calculation}

Suppose

Input Image

\[
32\times32\times3
\]

Convolution Layer

\begin{itemize}
\item Filters = 16
\item Kernel = $3\times3$
\end{itemize}

Parameters

\[
(3\times3\times3+1)\times16
\]

\[
=(27+1)\times16
\]

\[
=448
\]

Therefore,

Total trainable parameters = 448.

%%%%%%%%%%%%%%%%%%%%%%%%%%%%%%%%%%%%%%%%%%%%%%%%%%%%%%%%%%%%

\section*{7. Dataset}

Dataset: CIFAR-10

\begin{itemize}
\item Training Images : 50,000
\item Testing Images : 10,000
\item Classes : 10
\item Image Size : $32\times32\times3$
\end{itemize}

The dataset was loaded directly from the pickled batch files
(\texttt{data\_batch\_1} to \texttt{data\_batch\_5} for training and
\texttt{test\_batch} for testing) already stored in Google Drive from the
previous experiment, rather than downloaded again. Each batch file
unpickles to a dictionary whose \texttt{data} field is a
$10000 \times 3072$ array of pixel values and whose \texttt{labels}
field is a list of 10000 class indices; this was reshaped into
$10000 \times 32 \times 32 \times 3$ images before use.

%%%%%%%%%%%%%%%%%%%%%%%%%%%%%%%%%%%%%%%%%%%%%%%%%%%%%%%%%%%%

\section*{8. Experimental Procedure and Results}

\subsection*{Task 1: Dataset Preparation}

CIFAR-10 was loaded from Drive and the pixel values were normalised to
$[0,1]$, confirmed by a pixel range of $[0.00, 1.00]$. The training set
had 50,000 images and the test set had 10,000 images, each
$32\times32\times3$, across the 10 usual classes. Ten sample images are
shown in Figure~1, and the class distribution of the training set is
shown in Figure~2.

\begin{figure}[H]
\centering
\includegraphics[width=0.85\textwidth]{figures/01_sample_images}
\caption{Ten Sample CIFAR-10 Images}
\end{figure}

\begin{figure}[H]
\centering
\includegraphics[width=0.7\textwidth]{figures/02_class_distribution}
\caption{Class Distribution (Training Set)}
\end{figure}

\textbf{Observation:} The training set has exactly 5,000 images in each
of the 10 classes, so it is perfectly balanced and does not need any
class weighting or resampling before training.

\subsection*{Task 2: Convolution Layer — Kernel Size Comparison}

A single-filter convolution with `valid' padding and stride 1 was applied
to one sample image using 3x3, 5x5 and 7x7 kernels.

\begin{center}
\begin{tabular}{|c|c|c|}
\hline
Kernel Size & Output Shape & Expected Size (Formula)\\
\hline
$3\times3$ & $30\times30$ & 30\\
$5\times5$ & $28\times28$ & 28\\
$7\times7$ & $26\times26$ & 26\\
\hline
\end{tabular}
\end{center}

\begin{figure}[H]
\centering
\includegraphics[width=0.9\textwidth]{figures/03_kernel_size_comparison}
\caption{Feature Maps for Different Kernel Sizes}
\end{figure}

\textbf{Observation:} All three output sizes match the
$\frac{N-F+2P}{S}+1$ formula from Section 3 exactly, using $N=32$,
$P=0$, $S=1$. Larger kernels lose more pixels around the border, so the
output feature map shrinks as the kernel size increases.

\subsection*{Task 3: Stride and Padding}

The same demonstration image and a fixed 3x3 kernel were used to compare
stride 1 against stride 2, and `same' padding against `valid' padding.

\begin{center}
\begin{tabular}{|c|c|c|}
\hline
Stride & Padding & Output Shape\\
\hline
1 & Valid & $30\times30$\\
2 & Valid & $15\times15$\\
1 & Same & $32\times32$\\
2 & Same & $16\times16$\\
\hline
\end{tabular}
\end{center}

\textbf{Observation:} `Same' padding keeps the output the same size as
the input when the stride is 1, while `valid' padding always shrinks it.
Increasing the stride from 1 to 2 roughly halves the output size in both
padding modes, which again matches the formula from Section 3.

\subsection*{Task 4: Feature Map Visualization}

A freshly initialized (untrained) Conv2D layer with 16 filters was
applied to a sample image, and 8 of the resulting feature maps are shown
in Figure~3.

\begin{figure}[H]
\centering
\includegraphics[width=0.85\textwidth]{figures/04_feature_maps}
\caption{Feature Maps After the First Convolution Layer (Untrained)}
\end{figure}

\textbf{Observation:} Since the filters are randomly initialised and not
yet trained, most of them respond to roughly the same regions of the
image, mainly the brighter areas around the frog's body. Filter 5 stays
almost completely dark, which happens when a filter's random weights
combine with the ReLU activation in a way that zeroes out nearly its
entire output for this particular image -- a filter can end up doing
essentially nothing until training adjusts its weights.

\subsection*{Task 6: Building and Training the Main CNN}

The specified architecture, Input $\rightarrow$ Conv(16, 3x3)
$\rightarrow$ ReLU $\rightarrow$ MaxPool $\rightarrow$ Conv(32, 3x3)
$\rightarrow$ ReLU $\rightarrow$ MaxPool $\rightarrow$ Flatten
$\rightarrow$ Dense(10) $\rightarrow$ Softmax, was trained with the Adam
optimizer, 20 epochs and batch size 32, on 45,000 training images with
5,000 held out for validation, on the Google Colab GPU runtime. The
model had 25,578 parameters in total. Training took 115.95 seconds. By
the last epoch, training accuracy reached 76.61\% and validation accuracy
reached 66.40\%.

\begin{figure}[H]
\centering
\includegraphics[width=0.72\textwidth]{figures/05_training_accuracy}
\caption{Training Accuracy vs Epoch}
\end{figure}

\textbf{Observation:} Training accuracy rises quickly over the first few
epochs and then continues climbing more slowly, reaching 76.61\% by
epoch 20 without flattening out completely, which suggests a few more
epochs could still improve it slightly.

\begin{figure}[H]
\centering
\includegraphics[width=0.72\textwidth]{figures/06_validation_accuracy}
\caption{Validation Accuracy vs Epoch}
\end{figure}

\textbf{Observation:} Validation accuracy rises quickly at first, then
levels off around 66--67\% from about epoch 9 onward while training
accuracy keeps increasing. This growing gap between training and
validation accuracy (76.61\% against 66.40\% by the last epoch) is a sign
of mild overfitting, which is expected here since this small network has
no dropout or other regularization.

\begin{figure}[H]
\centering
\includegraphics[width=0.72\textwidth]{figures/07_training_loss}
\caption{Training Loss vs Epoch}
\end{figure}

\textbf{Observation:} Training loss decreases steadily and smoothly
across all 20 epochs, from about 1.53 to 0.68, with no bumps.

\begin{figure}[H]
\centering
\includegraphics[width=0.72\textwidth]{figures/08_validation_loss}
\caption{Validation Loss vs Epoch}
\end{figure}

\textbf{Observation:} Validation loss drops quickly for the first few
epochs, reaches its lowest point around epoch 9--10, and then rises
slowly for the rest of training even as training loss keeps falling --
the same overfitting pattern seen in the accuracy curves.

\subsection*{Task 7: Model Evaluation}

The trained model was evaluated on the full 10,000-image test set.

\begin{center}
\begin{tabular}{|p{6cm}|p{4cm}|}
\hline
Metric & Value\\
\hline
Testing Accuracy & 65.84\%\\
\hline
Weighted Precision & 66.77\%\\
\hline
Weighted Recall & 65.84\%\\
\hline
Weighted F1-score & 66.04\%\\
\hline
\end{tabular}
\end{center}

\begin{center}
\begin{tabular}{|l|c|c|c|c|}
\hline
Class & Precision & Recall & F1-score & Support\\
\hline
Airplane & 0.68 & 0.71 & 0.69 & 1000\\
Automobile & 0.83 & 0.74 & 0.78 & 1000\\
Bird & 0.60 & 0.50 & 0.55 & 1000\\
Cat & 0.43 & 0.53 & 0.47 & 1000\\
Deer & 0.59 & 0.62 & 0.60 & 1000\\
Dog & 0.58 & 0.55 & 0.57 & 1000\\
Frog & 0.81 & 0.66 & 0.73 & 1000\\
Horse & 0.67 & 0.75 & 0.71 & 1000\\
Ship & 0.72 & 0.81 & 0.76 & 1000\\
Truck & 0.77 & 0.72 & 0.74 & 1000\\
\hline
\textbf{Accuracy} & \multicolumn{3}{c|}{0.66} & 10000\\
\hline
Macro Avg & 0.67 & 0.66 & 0.66 & 10000\\
\hline
Weighted Avg & 0.67 & 0.66 & 0.66 & 10000\\
\hline
\end{tabular}
\end{center}

\begin{figure}[H]
\centering
\includegraphics[width=0.75\textwidth]{figures/09_confusion_matrix}
\caption{Confusion Matrix -- Main CNN (Max Pooling)}
\end{figure}

\textbf{Observation:} The largest confusion is between cat and dog: 171
cats were predicted as dog and 241 dogs were predicted as cat, matching
cat's precision of only 0.43, the lowest in the table. Automobile and
truck are also mixed up in both directions (123 automobiles predicted as
truck, 81 trucks predicted as automobile), and airplane and ship show
some mutual confusion as well (121 airplanes predicted as ship, 71 ships
predicted as airplane) -- in each case the confused classes share a
similar overall shape or background at this resolution.

%%%%%%%%%%%%%%%%%%%%%%%%%%%%%%%%%%%%%%%%%%%%%%%%%%%%%%%%%%%%

\section*{9. Mandatory Plots}

All of the plots required for this experiment -- sample images, class
distribution, training accuracy, validation accuracy, training loss,
validation loss, feature maps and the confusion matrix -- are presented
above in Section 8, alongside the task they belong to, each with its own
brief observation.

%%%%%%%%%%%%%%%%%%%%%%%%%%%%%%%%%%%%%%%%%%%%%%%%%%%%%%%%%%%%

\section*{10. Additional Exercises}

\subsection*{Exercise 1: Output Size for a 64x64 Image, 5x5 Kernel,
Stride 2, Padding 2}

Using the formula from Section 3 with $N=64$, $F=5$, $S=2$, $P=2$,

\[
\left\lfloor\frac{64-5+2(2)}{2}\right\rfloor+1=32
\]

This was checked against an actual convolution layer (zero-padded by 2,
then a $5\times5$ kernel with stride 2), which produced a $32\times32$
output, confirming the calculation.

\subsection*{Exercise 2: Trainable Parameters for 64 Filters, 3x3
Kernel, RGB Input}

\[
(3\times3\times3+1)\times64=1792
\]

This was also checked against an actual Conv2D layer with 64 filters, a
$3\times3$ kernel and a 3-channel input, which reported 1792 trainable
parameters, confirming the calculation.

\subsection*{Exercise 3: ReLU vs Sigmoid Activation}

The convolutional layers of the main CNN were switched from ReLU to
Sigmoid activation, keeping everything else identical, and trained for
the same 20 epochs.

\begin{center}
\begin{tabular}{|l|c|c|}
\hline
Activation & Test Accuracy & Training Time (s)\\
\hline
ReLU (main CNN) & 65.84\% & 115.95\\
Sigmoid & 60.04\% & 101.52\\
\hline
\end{tabular}
\end{center}

\textbf{Observation:} ReLU reached noticeably higher accuracy than
Sigmoid under the same settings (65.84\% against 60.04\%). By the last
epoch the Sigmoid network was still climbing steadily (training accuracy
60.93\%, well below ReLU's 76.61\%), which is consistent with Sigmoid
activations being more prone to small gradients that slow down learning,
especially in the earlier layers of the network.

\subsection*{Exercise 4: Replace Max Pooling with Average Pooling}

This is the same comparison as Task 5, so the two pooling types are
compared once, together, below.

\subsection*{Task 5 / Exercise 4: Max Pooling vs Average Pooling}

The main CNN already is the max-pooling version of this comparison. A
second copy of the same architecture, with `MaxPooling2D` replaced by
`AveragePooling2D` everywhere, was trained under identical settings.

\begin{center}
\begin{tabular}{|l|c|c|c|}
\hline
Pooling & Test Accuracy & Output Shape After 1st Pool & Training Time (s)\\
\hline
Max & 65.84\% & $16\times16$ & 115.95\\
Average & 67.38\% & $16\times16$ & 98.47\\
\hline
\end{tabular}
\end{center}

\begin{figure}[H]
\centering
\includegraphics[width=0.6\textwidth]{figures/10_pooling_comparison}
\caption{Max Pooling vs Average Pooling}
\end{figure}

\textbf{Observation:} Both pooling types reduce the spatial size in
exactly the same way, from $32\times32$ down to $16\times16$ after the
first pooling layer, since the output size only depends on the pooling
window and stride, not on whether the maximum or the average is taken.
Average pooling reached slightly higher test accuracy than max pooling
here (67.38\% against 65.84\%) and also trained a little faster. This is
not guaranteed to hold in general -- max pooling usually keeps the
strongest activations, which can help when the most distinctive feature
in a region matters most, while average pooling smooths over the whole
region, which can help when the useful information is spread out rather
than concentrated at a single point.

\subsection*{Exercise 5: Increasing Filters from 16 to 64}

The first convolution layer's filter count was increased from 16 to 64,
keeping the second layer at 32 filters, and trained under identical
settings.

\begin{center}
\begin{tabular}{|l|c|c|c|}
\hline
First-Layer Filters & Parameters & Test Accuracy & Training Time (s)\\
\hline
16 (main CNN) & 25,578 & 65.84\% & 115.95\\
64 & 40,746 & 67.05\% & 131.77\\
\hline
\end{tabular}
\end{center}

\begin{figure}[H]
\centering
\includegraphics[width=0.6\textwidth]{figures/11_filter_count_comparison}
\caption{Effect of Increasing Filters (16 to 64)}
\end{figure}

\textbf{Observation:} Increasing the first layer's filter count from 16
to 64 gave a small accuracy improvement (65.84\% to 67.05\%) at the cost
of more parameters (25,578 to 40,746) and a longer training time (115.95
seconds to 131.77 seconds). The accuracy gain here is modest compared to
the added computation, which suggests this particular network was not
strongly limited by the number of first-layer filters to begin with.

%%%%%%%%%%%%%%%%%%%%%%%%%%%%%%%%%%%%%%%%%%%%%%%%%%%%%%%%%%%%

\section*{11. Results}

\begin{center}
\begin{tabular}{|p{6cm}|p{4cm}|}
\hline
Metric & Value\\
\hline
Training Accuracy & 76.61\%\\
\hline
Testing Accuracy & 65.84\%\\
\hline
Precision & 66.77\%\\
\hline
Recall & 65.84\%\\
\hline
F1-score & 66.04\%\\
\hline
Number of Parameters & 25,578\\
\hline
\end{tabular}
\end{center}

%%%%%%%%%%%%%%%%%%%%%%%%%%%%%%%%%%%%%%%%%%%%%%%%%%%%%%%%%%%%

\section*{12. Discussion}

\begin{enumerate}

\item \textbf{Why is convolution preferred over fully connected layers
for images?}
A fully connected layer treats every pixel as an independent input and
learns a separate weight for each one, which for a $32\times32\times3$
image alone already means thousands of weights per neuron, and it throws
away the spatial relationship between neighbouring pixels entirely.
Convolution instead slides a small kernel across the image, so the same
small set of weights is reused at every position, giving far fewer
parameters and letting the network learn local patterns, such as edges
and textures, that are meaningful regardless of where they appear in the
image.

\item \textbf{How does stride affect the feature map size?}
Stride controls how many pixels the kernel moves between each
application. A stride of 1 slides the kernel one pixel at a time and
produces the largest possible output, while a larger stride skips more
pixels and shrinks the output size roughly in proportion, as shown
directly in Task 3, where changing the stride from 1 to 2 roughly halved
the output size in both padding modes.

\item \textbf{What is the role of padding?}
Padding adds extra pixels, usually zeros, around the border of the input
before convolution. Without it (`valid' padding), the output shrinks with
every convolution layer and pixels near the edges are used in fewer
computations than pixels near the centre. `Same' padding is chosen so
that the output has the same spatial size as the input, which was
confirmed directly in Task 3.

\item \textbf{Why is pooling used?}
Pooling reduces the spatial size of the feature maps, which cuts down the
number of parameters and computation needed in the layers that follow,
and it also makes the network's output a little less sensitive to small
shifts or distortions in exactly where a feature appears in the image.
Task 5 showed this reduction directly: both max and average pooling
brought the feature map down from $32\times32$ to $16\times16$ after the
first pooling layer.

\item \textbf{How do feature maps represent image characteristics?}
Each feature map is the result of one filter being applied across the
whole image, so it shows where and how strongly that particular pattern
(an edge at some orientation, a texture, a corner, and so on) is present.
Early layers tend to pick up simple, low-level patterns like edges and
colour contrasts, which is consistent with what Task 4 showed for an
untrained layer -- most filters responding to the same bright regions of
the sample image, since at that point they only reflect random initial
weights rather than anything learned yet.

\item \textbf{Why do CNNs require fewer parameters than MLPs?}
An MLP's first layer would need a separate weight for every pixel in the
image for every neuron in that layer, so the parameter count grows
directly with the image size. A CNN's convolution layer instead has one
small kernel per filter, reused at every spatial position across the
whole image, so the parameter count depends only on the kernel size and
the number of filters, not on the image's width or height. Task 6's main
CNN illustrates this directly: with two small convolution layers it has
only 25,578 parameters in total, most of which actually come from the
final dense layer rather than the convolutions themselves.

\end{enumerate}

%%%%%%%%%%%%%%%%%%%%%%%%%%%%%%%%%%%%%%%%%%%%%%%%%%%%%%%%%%%%

\section*{13. References}

\begin{enumerate}

\item Ian Goodfellow, Yoshua Bengio and Aaron Courville,
\emph{Deep Learning}, MIT Press, 2016.

\item Christopher M. Bishop,
\emph{Pattern Recognition and Machine Learning}, Springer, 2006.

\item Simon Haykin,
\emph{Neural Networks and Learning Machines}, Pearson, 2009.

\item TensorFlow Documentation:
\url{https://www.tensorflow.org}

\item Alex Krizhevsky, \emph{Learning Multiple Layers of Features from
Tiny Images} (CIFAR-10 dataset documentation):
\url{https://www.cs.toronto.edu/~kriz/cifar.html}

\end{enumerate}

%%%%%%%%%%%%%%%%%%%%%%%%%%%%%%%%%%%%%%%%%%%%%%%%%%%%%%%%%%%%

\end{document}
